\chapter{Taller Multi-Motor y Casos de Negocio (PostgreSQL / MySQL)}
\label{chap:lab02_taller_multimotor}

\section{Desarrollo Pr\'actico Multi-Motor y Compatibilidad}

A lo largo del dise\~no de bases de datos, los ingenieros deben dominar las particularidades sint\'acticas entre los dos motores de c\'odigo abierto m\'as utilizados del mundo: **PostgreSQL** y **MySQL**.

\subsection{Comparativa de Sintaxis y Dialectos}
\begin{table}[h!]
\centering
\small
\begin{tabularx}{\textwidth}{lXX}
\toprule
\textbf{Caracter\'istica} & \textbf{PostgreSQL 16+} & \textbf{MySQL 8.0+} \\
\midrule
Autoincrementables & \texttt{SERIAL} / \texttt{IDENTITY} & \texttt{AUTO\_INCREMENT} \\
Manejo de Conflictos & \texttt{ON CONFLICT ... DO UPDATE} & \texttt{ON DUPLICATE KEY UPDATE} \\
Retorno de inserciones & Cl\'ausula \texttt{RETURNING} & No soportado nativamente \\
Tipado Booleano & Nativo (\texttt{BOOLEAN}) & Alias de \texttt{TINYINT(1)} \\
JSON Avanzado & \texttt{JSONB} con \'indices GIN & \texttt{JSON} binario \\
\bottomrule
\end{tabularx}
\caption{Diferencias clave de sintaxis entre PostgreSQL y MySQL}
\end{table}

---

\section{Casos de Negocio Reales}

\subsection{Caso 1: Cineplex (Gesti\'on de Salas y Boletos)}
\begin{lstlisting}[style=sqlstyle, caption={Consulta de recaudaci\'on por sala de cine}]
SELECT
    c.cinema_name,
    ms.movie_title,
    COUNT(b.booking_id) AS tickets_sold,
    SUM(ms.ticket_price) AS gross_revenue
FROM cinemas AS c
INNER JOIN movie_showtimes AS ms
    ON c.cinema_id = ms.cinema_id
LEFT JOIN bookings AS b
    ON ms.showtime_id = b.showtime_id
GROUP BY
    c.cinema_name,
    ms.movie_title
ORDER BY
    gross_revenue DESC NULLS LAST;
\end{lstlisting}
